\documentclass{article}

% ----------------------------------------------------------------------
% 1. Package Inclusions & Submission Track Setup
% ----------------------------------------------------------------------
% Default option is for double-blind submission (adds line numbers).
% Remove [preprint] or [final] during the review phase.
\usepackage{neurips_2026}

\usepackage[utf8]{inputenc} % allow utf-8 input
\usepackage[T1]{fontenc}    % use 8-bit T1 fonts
\usepackage{hyperref}       % hyperlinks
\usepackage{url}            % simple URL typesetting
\usepackage{booktabs}       % professional-quality tables
\usepackage{amsfonts}       % blackboard math symbols
\usepackage{amsmath}        % advanced math typesetting (highly recommended)
\usepackage{amssymb}        % additional math symbols
\usepackage{nicefrac}       % compact symbols for 1/2, etc.
\usepackage{microtype}      % microtypography
\usepackage{xcolor}         % colors
\usepackage{graphicx}       % required for including figures

% ----------------------------------------------------------------------
% 2. Paper Title & Author Information
% ----------------------------------------------------------------------
\title{CORE: Centered Orthogonal Routing for Expert Specialization in Mixture-of-Experts}
% The \author macro works with any number of authors.
% Keep it anonymous for the initial submission (leave as placeholder).
\author{%
  Anonymous Author(s)\\
  Affiliation\\
  Address\\
  \texttt{email} \\
}

\begin{document}

\maketitle

% ----------------------------------------------------------------------
% 3. Abstract
% ----------------------------------------------------------------------
\begin{abstract}
Mixture-of-Experts (MoE) scales model capacity only when routing induces meaningful expert specialization. We therefore study routing through the lens of specialization rather than treating continual learning as the primary theoretical object. Our starting observation is that current routing pipelines can improve occupancy statistics without necessarily inducing clean specialization. We trace this gap to the anisotropic geometry of language-model representations: a large global mean shift makes uncentered routing collapse-prone, yet the mean direction itself still carries useful coarse routing information. This motivates a centered and orthogonally decoupled first-order gate that explicitly separates mean-direction routing from centered residual routing. We instantiate this idea with CORE. Empirically, we show that real token representations are strongly mean-separable after centering, that CORE achieves the highest routing separation ratio among the compared gates, that function words are stably absorbed by a small set of experts under the decoupled routing mechanism, and that adding second-order complexity brings substantial optimization competition and engineering overhead without corresponding geometric gains. These results support a simple conclusion: in the current regime, expert specialization is best served by centered first-order orthogonal decoupling, while continual learning is better treated as a downstream validation setting.
\end{abstract}

% ----------------------------------------------------------------------
% 4. Main Content (Imported from sections/)
% ----------------------------------------------------------------------
% We use \input{} to keep the main file clean and modular.

\section{Introduction}
\label{sec:intro}

Mixture-of-Experts (MoE) improves nominal model capacity by activating only a small subset of parameters for each token, but the scientific foundation of MoE is not sparsity by itself. It is \emph{expert specialization}. An MoE becomes genuinely useful only when different experts learn meaningfully different transformations, respond to meaningfully different parts of the data distribution, and together convert architectural multiplicity into \emph{effective} capacity. If this does not happen, additional experts remain largely redundant modules: the model pays the systems cost of multiple experts without obtaining the intended functional diversity \citep{fedus2022switch}. From this viewpoint, expert specialization is not an auxiliary property of MoE; it is the condition under which MoE works at all.

In practice, however, expert specialization is often much weaker than the architecture seems to suggest. Experts can look balanced in usage yet remain functionally similar; they can appear dispersed in parameter space yet still absorb highly correlated learning signals; and they can process superficially different tokens while relying on the same underlying representational directions. \textcolor{red}{Here we should insert the finalized empirical evidence from our main-model audit and probe experiments, showing exactly where current experts fail to specialize and which commonly used proxies break down.} The important point is that current failure is not exhausted by occupancy imbalance, expert under-utilization, or a single overlap statistic. The deeper problem is that the system often fails to produce experts that are simultaneously different in parameters, different in the data manifolds they process, and different in the updates they accumulate.

This also clarifies the limitation of many existing lines of work. Standard MoE training specifies specialization mostly \emph{indirectly}: sparse activation decides which experts are active, load-balancing objectives regulate how often each expert is used, shared-expert or fallback designs absorb common knowledge, and various regularizers encourage experts to look different in parameter space. These mechanisms can improve utilization, stability, or apparent diversity, but they do not directly enforce the stronger condition we actually care about: that the gate and the experts cooperate to form a system in which different experts consistently absorb different semantic submanifolds and therefore evolve along different learning trajectories. Parameter independence by itself is not enough; occupancy balance by itself is not enough; even geometric separation by itself is not enough. A more essential notion of specialization must jointly involve \emph{which data manifold is routed where} and \emph{how that routing reshapes expert-specific updates}.

Our current analysis is anchored in the geometry of hidden representations, especially mean bias and anisotropy \citep{ethayarajh2019contextual,chi2022representation}. Language-model activations contain a strong shared/common component, and this common structure can dominate similarity relationships before any expert-specific signal becomes visible. As a result, tokens that should be separated for specialization may first appear close along shared, high-energy directions, while semantically discriminative structure is relegated to weaker residual components. This makes representational geometry a natural starting point: if the input to the expert system is already dominated by a common mode, then apparent diversity at the architectural level need not translate into true specialization at the learning level.

The key lesson of our current milestone is therefore negative but important: making representations look more separable does not automatically make experts more specialized. Our recent experiments suggest that the missing link is the coupling between representational separation and expert-conditioned gradient differentiation. \textcolor{red}{This is where the final experimental section should show, with finalized metrics and ablations, that cleaner routing geometry or better-looking separation does not automatically imply stronger expert specialization under a stricter criterion.} This reframes the main question of the paper. Rather than starting from a fixed routing design and asking whether it works, we ask a more foundational question: \emph{what should count as expert specialization, and what kind of gate--expert system can actually realize it?}

Accordingly, the goal of this work is to move from superficial proxies of specialization toward a more intrinsic measure. In the current draft, we treat expert specialization as a joint property of three coupled objects: expert parameters, the routed data submanifolds, and the expert-conditioned updates induced by those routed data. This viewpoint is directly motivated by our empirical pipeline and will guide the remainder of the paper. \textcolor{red}{The final method is intentionally left open at this stage. The present draft should first establish the problem formulation, the failure diagnosis, and the evaluation target; a definitive routing architecture should be introduced only after the corresponding evidence becomes mature.}

At the current checkpoint, the introduction is organized around the following claims:
\begin{itemize}
    \item \textbf{Expert specialization is a coupled systems property.} It cannot be reduced to parameter overlap, expert load, or any single occupancy statistic.
    \item \textbf{A stronger notion of specialization must connect experts to data manifolds and updates.} Experts are specialized only when the routed data they receive and the gradients they accumulate both become differentiated.
    \item \textbf{Mean bias and anisotropy provide the current diagnostic anchor.} The representational common mode is a plausible upstream source of weak specialization, but \textcolor{red}{the final operational metric and the final routing mechanism remain to be finalized by the next round of experiments.}
\end{itemize}

In this sense, the contribution of the current draft is to sharpen the problem itself: what MoE needs is not merely a more elaborate expert inventory, but a principled account of how specialization is created, measured, and sustained.


\input{sections/2related.tex}

\section{Understanding Expert Specialization}
\label{sec:analysis}

Our analysis starts from a narrower question than continual learning: what kind of routing geometry actually induces expert specialization in modern MoE language models? The answer we develop in this section is first-order and decoupled. In the current regime, the key is not to add more nonlinear complexity, but to isolate two distinct first-order signals that standard gates entangle: the dominant mean-related direction and the centered residual variation.

\subsection{Setup and Notation}
Let $x \in \mathbb{R}^d$ denote the hidden representation before routing. Let $\mu_{\mathrm{global}}$ be the global mean. Define
$$
u=\frac{\mu_{\mathrm{global}}}{\|\mu_{\mathrm{global}}\|_2}, \qquad
x_{\parallel}=(u^\top x)u, \qquad
x_c=x-\mu_{\mathrm{global}}.
$$

Here, $x_{\parallel}$ captures the coordinate of the token along the dominant mean-related direction, while $x_c$ removes the shared global offset and retains the centered residual variation.

\subsection{Why Centering Is Necessary}
The correct statement is not that centering makes conditional means equal. It does the opposite: it removes the shared global offset and preserves the domain-specific first-order differences that were previously buried under that offset.

\begin{proposition}[Centering preserves domain-specific first-order differences]
Let
$$
\mu_A=\mathbb{E}[x\mid A], \qquad \mu_B=\mathbb{E}[x\mid B].
$$
Then
$$
\mathbb{E}[x_c\mid A]=\mu_A-\mu_{\mathrm{global}}, \qquad
\mathbb{E}[x_c\mid B]=\mu_B-\mu_{\mathrm{global}}.
$$
Hence
$$
\mathbb{E}[x_c\mid A]-\mathbb{E}[x_c\mid B]=\mu_A-\mu_B.
$$
Therefore centering removes only the shared offset; it does not erase domain-specific first-order separation.
\end{proposition}

\begin{proof}
The identities follow directly from linearity of conditional expectation:
$$
\mathbb{E}[x_c\mid A]
=
\mathbb{E}[x-\mu_{\mathrm{global}}\mid A]
=
\mathbb{E}[x\mid A]-\mu_{\mathrm{global}}
=
\mu_A-\mu_{\mathrm{global}},
$$
and similarly for $B$. Subtracting the two equalities yields the claim.
\end{proof}

This proposition explains why centered linear routing can already become much cleaner than uncentered routing. Uncentered routing compares tokens through a large shared bias term. Centered routing exposes the differences that matter for assignment.

\subsection{Why the Mean Direction Must Be Kept}
Centering is necessary, but it is not sufficient. The mean direction cannot simply be discarded, because it still carries coarse routing information. In practice, the dominant direction often correlates with token frequency, function-word behavior, or broad domain drift. If the gate removes that direction completely, it loses a stable and useful routing cue.

This leads to a structured first-order interpretation:
$$
\ell_e(x)=a_e^\top x_{\parallel}+b_e^\top x_c.
$$
The first term handles coarse, mean-related separation; the second term handles centered residual separation.

\begin{proposition}[Orthogonally decoupled first-order routing]
Define
$$
\Delta_{\parallel}
=
\mathbb{E}[x_{\parallel}\mid A]-\mathbb{E}[x_{\parallel}\mid B],
\qquad
\Delta_c
=
\mathbb{E}[x_c\mid A]-\mathbb{E}[x_c\mid B].
$$
Then
$$
\mathbb{E}[\ell_e(x)\mid A]-\mathbb{E}[\ell_e(x)\mid B]
=
a_e^\top \Delta_{\parallel}+b_e^\top \Delta_c.
$$
Therefore the two channels can independently exploit mean-direction differences and centered first-order differences.
\end{proposition}

\begin{proof}
Take conditional expectations under $A$ and $B$ and subtract linearly.
\end{proof}

This is the theoretical basis of CORE. The gate is not stronger because it is more nonlinear. It is stronger because it exposes the right first-order factors explicitly.

\subsection{Specialization and Low-Rankness}
Expert specialization is an assignment-side phenomenon: similar tokens are routed to the same expert more consistently. Low-rankness is a geometry-side phenomenon: activations arriving at one expert occupy a smaller effective subspace. The two are related through the law of total covariance:
$$
\operatorname{Cov}(x_c)
=
\mathbb{E}[\operatorname{Cov}(x_c\mid E)]
+
\operatorname{Cov}(\mathbb{E}[x_c\mid E]),
$$
where $E$ is the routed expert.

This decomposition tells us how to read specialization geometrically. The second term measures how far the expert-conditional means are pulled apart. The first term measures how spread out tokens remain inside each expert. If routing increases the second term under a fixed global covariance budget, then the first term is pressured downward. In anisotropic language-model representations, this pressure naturally encourages lower effective within-expert rank.

We therefore use two complementary metrics:
$$
\rho_{\mathrm{sep}}
=
\frac{\operatorname{Tr}(\operatorname{Cov}(\mathbb{E}[x_c\mid E]))}
{\operatorname{Tr}(\operatorname{Cov}(x_c))}
$$
and
$$
\mathrm{sRank}(\Sigma_e)
=
\frac{\operatorname{Tr}(\Sigma_e)^2}{\operatorname{Tr}(\Sigma_e^2)},
\qquad
\Sigma_e=\operatorname{Cov}(x_c\mid E=e).
$$

$\rho_{\mathrm{sep}}$ measures how strongly the gate separates experts in mean space. $\mathrm{sRank}(\Sigma_e)$ measures whether each expert's assigned activations collapse toward a lower-dimensional subspace. Specialization is strongest when these two move in the same direction.

\subsection{Current Empirical Evidence}
The current analysis already supports four concrete claims.

\paragraph{Data geometry is predominantly mean-separable.}
We extracted centered token activations from six representative MoE layers and compared K-Means clustering against subspace-based clustering. Across all six layers, K-Means produced consistently higher separation ratios than K-Subspaces. This indicates that real token clouds are primarily separated by location in residual space, not by hidden covariance-only structure. In other words, after centering, the representation geometry still looks more like separated clusters than like coincident ellipsoids distinguished only by shape.

\paragraph{The decoupled first-order gate achieves the strongest separation.}
On fixed checkpoint audits, CORE achieves the highest mean separation ratio among the compared routing families:
$$
\rho_{\mathrm{sep}}(\texttt{CORE})=0.1015,
$$
which exceeds `hybrid' $(0.0645)$, `linear\_running\_center' $(0.0615)$, and `linear\_no\_center' $(0.0060)$. This is the clearest direct evidence that expert specialization in the current regime is governed by mean separation, not by higher-order slicing.

\paragraph{The decoupling has a real semantic interpretation.}
Using a strict single-token vocabulary of 43 function words and 109 domain keywords, we observe a strong function-word sink effect. The average top-share for function words exceeds $0.60$ under CORE and the hybrid comparator, compared with $0.5778$ under the ordinary centered linear baseline. This supports the interpretation that the mean-related channel absorbs high-frequency or function-like tokens, leaving the residual channel to handle finer semantic distinctions.

\paragraph{Second-order complexity is not justified in the current regime.}
The second-order branch was not numerically dead: during early training, its gradient norm was of the same order as the centered first-order branch, with average ratio $0.8685$. This matters because it rules out the simplest failure explanation, namely that the higher-order term is merely starved by optimization. At the same time, the engineering cost is prohibitive. The hybrid router requires roughly $7.23$M router parameters and $1.196$M FLOPs/token, while CORE uses only $221.2$K parameters and $15.4$K FLOPs/token. Since the geometry audit already shows that the current data regime is predominantly mean-separable, the extra complexity brings no commensurate specialization benefit. In this regime, the correct conclusion is not that second-order routing is impossible, but that it is unnecessary.

\subsection{Planned Evidence That Strengthens the Story}
Two additional analyses are planned to make the story fully publication-ready.

\paragraph{Eigen-spectrum decay audit.}
We will compare the local covariance spectra of experts under CORE and under the no-centering baseline. If our current interpretation is correct, CORE should produce a visibly sharper eigenvalue decay, showing that stronger mean separation passively induces lower effective within-expert dimensionality.

\paragraph{Orthogonal ablation and continual-learning bridge.}
We will compare three pure first-order variants: mean-only routing, residual-only routing, and the full decoupled gate. This will test whether the gain truly comes from orthogonal decoupling rather than from either channel alone. We then use continual-learning interference as a downstream bridge, not as the primary theoretical object.

\subsection{Main Takeaway}
The current evidence supports a simple conclusion. The bottleneck in current MoE routing is not the absence of a stronger nonlinear gate. It is the failure to handle anisotropic first-order geometry correctly. Once the global offset is removed and the mean direction is preserved as a separate routing channel, expert specialization becomes substantially cleaner. This is why the analysis section supports a strong practical statement: for the current representation regime, second-order routing is not needed to explain or improve specialization.


\section{Centered Orthogonal Routing for Experts}
\label{sec:method}

The analysis above leads to a simple design principle: a useful MoE gate must remove the shared global offset, but it must not discard the mean-direction information that still carries coarse routing cues. CORE implements this idea with an explicitly decoupled first-order gate.

\subsection{Routing Decomposition}
Given the hidden representation $x$, we define
$$
u=\frac{\mu_{\mathrm{global}}}{\|\mu_{\mathrm{global}}\|_2}, \qquad
x_{\parallel}=(u^\top x)u, \qquad
x_c=x-\mu_{\mathrm{global}}.
$$
The router for expert $e$ is then
$$
\ell_e(x)=a_e^\top x_{\parallel}+b_e^\top x_c,
$$
where $a_e,b_e \in \mathbb{R}^d$.

This decomposition creates two explicit routing channels:
\begin{itemize}
    \item the \textbf{mean channel} $a_e^\top x_{\parallel}$, which captures coarse variation aligned with the dominant mean-related direction;
    \item the \textbf{residual channel} $b_e^\top x_c$, which captures centered first-order variation after the shared global offset has been removed.
\end{itemize}

\subsection{Why This Form Is Necessary}
The design is motivated by two complementary failure modes.

First, uncentered routing entangles a large shared bias with content-dependent variation. This makes the gate sensitive to the dominant global direction in a way that can produce collapse-like assignments.

Second, centered routing alone is incomplete if it simply discards the mean direction. The mean-related axis still contains useful routing information, especially for function words, high-frequency tokens, and other coarse token classes. CORE therefore preserves that axis explicitly, rather than forcing a single matrix to absorb both coarse and fine routing roles.

\subsection{Specialization-Oriented Objective}
CORE is trained with the standard task loss. In the current first-order version, we do not rely on quadratic routing terms. The method claim is therefore minimal: if the representation regime is predominantly mean-separable, then expert specialization can be improved by routing with a centered and orthogonally decoupled first-order gate.

\subsection{Task-Agnostic Routing}
CORE does not use task identity at inference time. Routing is computed directly from the current hidden representation. This makes the method naturally compatible with task-agnostic continual-learning protocols, although the value of this property must still be validated in downstream experiments.

\subsection{Computational Simplicity}
The computational cost of CORE is linear in the hidden dimension per expert. Both routing channels are first-order projections, so the overhead is modest and easy to analyze. This simplicity is not just an engineering convenience; it is part of the method claim. If the current regime is already well served by decoupled first-order routing, then introducing higher-order complexity is unnecessary.


\section{Experiments}
\label{sec:experiments}

Our experiments follow a single logic chain:
$$
\text{data geometry}
\rightarrow
\text{routing separation}
\rightarrow
\text{semantic decoupling}
\rightarrow
\text{rejection of unnecessary second-order complexity}.
$$
The goal is not to overwhelm the reader with many loosely connected diagnostics. The goal is to show, step by step, why expert specialization in the current regime is best explained by centered orthogonal first-order routing.

\subsection{Experiment 1: Data Geometry Audit}
\paragraph{Objective.}
We first ask a foundational question: after removing the global offset, are real token representations primarily separated by mean location or by covariance shape? This determines the upper bound on what a router should try to exploit.

\paragraph{Setup.}
We use a pretrained MoE model without imposing a special routing bias. From six representative MoE layers (layers 1, 3, 5, 7, 9, and 11), we extract 12{,}000 centered token activations $x_c$ from the WoS corpus.

\paragraph{Procedure.}
At each layer, we apply two offline clustering procedures:
\begin{itemize}
    \item K-Means, which searches for cluster separation through Euclidean distance and therefore emphasizes mean differences;
    \item K-Subspaces, which searches for low-rank subspace structure and therefore emphasizes covariance differences.
\end{itemize}

\paragraph{Metric.}
We quantify the geometric separability of the resulting partitions using the separation ratio
$$
\rho_{\mathrm{sep}}
=
\frac{\operatorname{Tr}(\operatorname{Cov}(\mathbb{E}[x_c\mid E]))}
{\operatorname{Tr}(\operatorname{Cov}(x_c))}.
$$

\paragraph{Result.}
Across all six layers, K-Means yields consistently higher $\rho_{\mathrm{sep}}$ than K-Subspaces. Moreover, the subspace-based partitions do not reliably produce lower within-cluster effective rank.

\paragraph{Conclusion.}
The residual representation geometry is predominantly mean-separable. Real token clouds behave more like separated groups at different locations than like clusters sharing a center but differing mainly by covariance shape. This result establishes the physical basis for using first-order hyperplanes as the primary routing mechanism.

\subsection{Experiment 2: Routing Separation Evaluation}
\paragraph{Objective.}
We next test whether the proposed routing mechanism actually achieves stronger expert separation than competing gates on real checkpoints.

\paragraph{Setup.}
We evaluate checkpoints at step 60{,}000 for four models:
\begin{itemize}
    \item `linear\_no\_center',
    \item `linear\_running\_center',
    \item `acp\_lite',
    \item `hybrid'.
\end{itemize}
We use the WoS validation split with 1{,}200 samples.

\paragraph{Procedure.}
For each model, we run forward inference, extract centered activations $x_c$ and expert assignments across all MoE layers, and aggregate geometry metrics layer by layer.

\paragraph{Metrics.}
We use two metrics:
$$
\rho_{\mathrm{sep}}
=
\frac{\operatorname{Tr}(\operatorname{Cov}(\mathbb{E}[x_c\mid E]))}
{\operatorname{Tr}(\operatorname{Cov}(x_c))}
$$
and
$$
\mathrm{sRank}(\Sigma_e)
=
\frac{\operatorname{Tr}(\Sigma_e)^2}{\operatorname{Tr}(\Sigma_e^2)}.
$$

\paragraph{Result.}
CORE (`acp\_lite') achieves the highest mean separation ratio:
$$
0.1015,
$$
compared with
$$
0.0645 \text{ for hybrid}, \qquad
0.0615 \text{ for centered linear}, \qquad
0.0060 \text{ for no-center linear}.
$$
At the same time, CORE maintains a reasonable within-expert effective-rank profile.

\paragraph{Conclusion.}
The strongest routing geometry in the current regime is achieved by orthogonally decoupled first-order routing. Once the dominant mean-related direction is separated from the centered residual channel, the gate maximizes between-expert mean separation without resorting to higher-order complexity.

\subsection{Experiment 3: Semantic Decoupling Audit}
\paragraph{Objective.}
The geometric decomposition must also correspond to a physically meaningful semantic decomposition. We therefore test whether the mean-direction channel and residual channel react to different token classes.

\paragraph{Setup.}
We use `linear\_running\_center', `acp\_lite', and `hybrid'. We construct a strict single-token vocabulary containing 43 function words and 109 domain keywords.

\paragraph{Procedure.}
We measure token assignment concentration through top-share statistics and, for the hybrid gate, decompose the pre-softmax logits into three parts:
$$
L_{\mathrm{mean}}=a_e^\top x_{\parallel}, \qquad
L_{\mathrm{residual}}=b_e^\top x_c, \qquad
L_{\mathrm{subspace}}=\|V_e x_c\|_2^2.
$$

\paragraph{Result.}
Function words exhibit a clear sink effect. Under `acp\_lite' and `hybrid', the mean top-share of function words exceeds $0.60$:
$$
0.6036 \text{ and } 0.6094,
$$
while the ordinary centered linear baseline reaches
$$
0.5778.
$$

\paragraph{Conclusion.}
The mean-direction channel carries real semantic workload. It acts as a stable absorber for high-frequency or function-like tokens, allowing the centered residual channel to focus on finer semantic separation. This makes the decoupling physically interpretable, not merely algebraically convenient.

\subsection{Experiment 4: Dynamics and Overhead Audit}
\paragraph{Objective.}
If first-order decoupling is already sufficient, why not keep second-order terms as an additional refinement? We answer this through optimization dynamics and engineering overhead.

\paragraph{Setup.}
For dynamics, we initialize the hybrid model and run a 100-step light training run. For cost, we evaluate stable checkpoints under fixed batch and token settings.

\paragraph{Procedure.}
We record the gradient $L_2$ norms of the pre-softmax components during early training and measure parameter count, router FLOPs per token, and wall-clock latency.

\paragraph{Metrics.}
For dynamics, we use
$$
\mathrm{Ratio}
=
\frac{\|\nabla L_{\mathrm{subspace}}\|_2}{\|\nabla L_{\mathrm{residual}}\|_2}.
$$
For cost, we report router parameters and FLOPs/token.

\paragraph{Result.}
The average gradient ratio is
$$
0.8685,
$$
which shows that the second-order branch is not trivially starved. By the final recorded step, the centered first-order and second-order branches still have comparable gradient magnitudes:
$$
0.00113 \text{ versus } 0.00201.
$$
At the same time, the engineering cost is dramatic. The hybrid router uses about $7.23$M router parameters and $1.196$M FLOPs/token, while CORE requires only $221.2$K router parameters and $15.4$K FLOPs/token.

\paragraph{Conclusion.}
In the current regime, second-order complexity is unnecessary. It is not absent because optimization starves it; rather, it competes with the first-order branch while offering no corresponding geometric gain, and it does so at nearly two orders of magnitude higher cost.

\subsection{Downstream View}
The four experiments above establish the main story of the paper. They show what structure exists in the data, what the router should exploit, why the proposed decoupling works, and why extra second-order complexity is not justified in the current regime. Continual learning is then used as a downstream validation setting: once specialization has been made geometrically cleaner, we can test whether this improved routing structure transfers into lower interference and better sequential adaptation.
 % Assuming you want experiments before related work based on usual flow, adjust if needed.

\section{Conclusion}
\label{sec:conclusion}

This paper reframes the routing problem in MoE around expert specialization. The central question is not whether a more complex gate can always be designed, but whether the gate separates tokens into experts in a geometrically meaningful way. Our analysis shows that the key obstacle in the current regime is anisotropic first-order geometry: a large global mean shift must be removed, but the mean direction itself still carries useful routing information.

This leads to a simple conclusion. The strongest routing mechanism in the current regime is a centered and orthogonally decoupled first-order gate. By separating mean-direction routing from centered residual routing, CORE induces cleaner expert specialization than standard linear gates, admits a direct semantic interpretation, and avoids the heavy cost of second-order complexity.

The current evidence is already sufficient to support this story through data-geometry audits, routing-separation evaluations, semantic decoupling analysis, and dynamics/overhead audits. The remaining work is downstream: we will continue to test whether the resulting expert specialization transfers into lower interference and improved continual-learning behavior.


% ----------------------------------------------------------------------
% 5. Acknowledgments (Hidden during double-blind review)
% ----------------------------------------------------------------------
\begin{ack}
  % Use unnumbered first level headings for the acknowledgments.
  % The ack environment automatically hides this section in the anonymized submission.
  % Declare funding and competing interests here.
\end{ack}

% ----------------------------------------------------------------------
% 6. References
% ----------------------------------------------------------------------
% The natbib package is loaded automatically. 
% We use your local 'refs.bib' file located in the root directory.
\bibliographystyle{plainnat}
\bibliography{refs}

% ----------------------------------------------------------------------
% 7. Appendix & Checklist
% ----------------------------------------------------------------------
\newpage
\appendix
% Technical appendices can be placed here directly or imported.
% \input{sections/appendix.tex}

\newpage
% The checklist MUST be included at the very end. Do not modify its structure.
\section*{NeurIPS Paper Checklist}

%%% BEGIN INSTRUCTIONS %%%
The checklist is designed to encourage best practices for responsible machine learning research, addressing issues of reproducibility, transparency, research ethics, and societal impact. Do not remove the checklist: {\bf The papers not including the checklist will be desk rejected.} The checklist should follow the references and follow the (optional) supplemental material.  The checklist does NOT count towards the page
limit. 

Please read the checklist guidelines carefully for information on how to answer these questions. For each question in the checklist:
\begin{itemize}
    \item You should answer \answerYes{}, \answerNo{}, or \answerNA{}.
    \item \answerNA{} means either that the question is Not Applicable for that particular paper or the relevant information is Not Available.
    \item Please provide a short (1--2 sentence) justification right after your answer (even for \answerNA). 
   % \item {\bf The papers not including the checklist will be desk rejected.}
\end{itemize}

{\bf The checklist answers are an integral part of your paper submission.} They are visible to the reviewers, area chairs, senior area chairs, and ethics reviewers. You will also be asked to include it (after eventual revisions) with the final version of your paper, and its final version will be published with the paper.

The reviewers of your paper will be asked to use the checklist as one of the factors in their evaluation. While \answerYes{} is generally preferable to \answerNo{}, it is perfectly acceptable to answer \answerNo{} provided a proper justification is given (e.g., error bars are not reported because it would be too computationally expensive'' or ``we were unable to find the license for the dataset we used''). In general, answering \answerNo{} or \answerNA{} is not grounds for rejection. While the questions are phrased in a binary way, we acknowledge that the true answer is often more nuanced, so please just use your best judgment and write a justification to elaborate. All supporting evidence can appear either in the main paper or the supplemental material, provided in appendix. If you answer \answerYes{} to a question, in the justification please point to the section(s) where related material for the question can be found.

IMPORTANT, please:
\begin{itemize}
    \item {\bf Delete this instruction block, but keep the section heading ``NeurIPS Paper Checklist"},
    \item  {\bf Keep the checklist subsection headings, questions/answers and guidelines below.}
    \item {\bf Do not modify the questions and only use the provided macros for your answers}.
\end{itemize} 
 

%%% END INSTRUCTIONS %%%


\begin{enumerate}

\item {\bf Claims}
    \item[] Question: Do the main claims made in the abstract and introduction accurately reflect the paper's contributions and scope?
    \item[] Answer: \answerTODO{} % Replace by \answerYes{}, \answerNo{}, or \answerNA{}.
    \item[] Justification: \justificationTODO{}
    \item[] Guidelines:
    \begin{itemize}
        \item The answer \answerNA{} means that the abstract and introduction do not include the claims made in the paper.
        \item The abstract and/or introduction should clearly state the claims made, including the contributions made in the paper and important assumptions and limitations. A \answerNo{} or \answerNA{} answer to this question will not be perceived well by the reviewers. 
        \item The claims made should match theoretical and experimental results, and reflect how much the results can be expected to generalize to other settings. 
        \item It is fine to include aspirational goals as motivation as long as it is clear that these goals are not attained by the paper. 
    \end{itemize}

\item {\bf Limitations}
    \item[] Question: Does the paper discuss the limitations of the work performed by the authors?
    \item[] Answer: \answerTODO{} % Replace by \answerYes{}, \answerNo{}, or \answerNA{}.
    \item[] Justification: \justificationTODO{}
    \item[] Guidelines:
    \begin{itemize}
        \item The answer \answerNA{} means that the paper has no limitation while the answer \answerNo{} means that the paper has limitations, but those are not discussed in the paper. 
        \item The authors are encouraged to create a separate ``Limitations'' section in their paper.
        \item The paper should point out any strong assumptions and how robust the results are to violations of these assumptions (e.g., independence assumptions, noiseless settings, model well-specification, asymptotic approximations only holding locally). The authors should reflect on how these assumptions might be violated in practice and what the implications would be.
        \item The authors should reflect on the scope of the claims made, e.g., if the approach was only tested on a few datasets or with a few runs. In general, empirical results often depend on implicit assumptions, which should be articulated.
        \item The authors should reflect on the factors that influence the performance of the approach. For example, a facial recognition algorithm may perform poorly when image resolution is low or images are taken in low lighting. Or a speech-to-text system might not be used reliably to provide closed captions for online lectures because it fails to handle technical jargon.
        \item The authors should discuss the computational efficiency of the proposed algorithms and how they scale with dataset size.
        \item If applicable, the authors should discuss possible limitations of their approach to address problems of privacy and fairness.
        \item While the authors might fear that complete honesty about limitations might be used by reviewers as grounds for rejection, a worse outcome might be that reviewers discover limitations that aren't acknowledged in the paper. The authors should use their best judgment and recognize that individual actions in favor of transparency play an important role in developing norms that preserve the integrity of the community. Reviewers will be specifically instructed to not penalize honesty concerning limitations.
    \end{itemize}

\item {\bf Theory assumptions and proofs}
    \item[] Question: For each theoretical result, does the paper provide the full set of assumptions and a complete (and correct) proof?
    \item[] Answer: \answerTODO{} % Replace by \answerYes{}, \answerNo{}, or \answerNA{}.
    \item[] Justification: \justificationTODO{}
    \item[] Guidelines:
    \begin{itemize}
        \item The answer \answerNA{} means that the paper does not include theoretical results. 
        \item All the theorems, formulas, and proofs in the paper should be numbered and cross-referenced.
        \item All assumptions should be clearly stated or referenced in the statement of any theorems.
        \item The proofs can either appear in the main paper or the supplemental material, but if they appear in the supplemental material, the authors are encouraged to provide a short proof sketch to provide intuition. 
        \item Inversely, any informal proof provided in the core of the paper should be complemented by formal proofs provided in appendix or supplemental material.
        \item Theorems and Lemmas that the proof relies upon should be properly referenced. 
    \end{itemize}

    \item {\bf Experimental result reproducibility}
    \item[] Question: Does the paper fully disclose all the information needed to reproduce the main experimental results of the paper to the extent that it affects the main claims and/or conclusions of the paper (regardless of whether the code and data are provided or not)?
    \item[] Answer: \answerTODO{} % Replace by \answerYes{}, \answerNo{}, or \answerNA{}.
    \item[] Justification: \justificationTODO{}
    \item[] Guidelines:
    \begin{itemize}
        \item The answer \answerNA{} means that the paper does not include experiments.
        \item If the paper includes experiments, a \answerNo{} answer to this question will not be perceived well by the reviewers: Making the paper reproducible is important, regardless of whether the code and data are provided or not.
        \item If the contribution is a dataset and\slash or model, the authors should describe the steps taken to make their results reproducible or verifiable. 
        \item Depending on the contribution, reproducibility can be accomplished in various ways. For example, if the contribution is a novel architecture, describing the architecture fully might suffice, or if the contribution is a specific model and empirical evaluation, it may be necessary to either make it possible for others to replicate the model with the same dataset, or provide access to the model. In general. releasing code and data is often one good way to accomplish this, but reproducibility can also be provided via detailed instructions for how to replicate the results, access to a hosted model (e.g., in the case of a large language model), releasing of a model checkpoint, or other means that are appropriate to the research performed.
        \item While NeurIPS does not require releasing code, the conference does require all submissions to provide some reasonable avenue for reproducibility, which may depend on the nature of the contribution. For example
        \begin{enumerate}
            \item If the contribution is primarily a new algorithm, the paper should make it clear how to reproduce that algorithm.
            \item If the contribution is primarily a new model architecture, the paper should describe the architecture clearly and fully.
            \item If the contribution is a new model (e.g., a large language model), then there should either be a way to access this model for reproducing the results or a way to reproduce the model (e.g., with an open-source dataset or instructions for how to construct the dataset).
            \item We recognize that reproducibility may be tricky in some cases, in which case authors are welcome to describe the particular way they provide for reproducibility. In the case of closed-source models, it may be that access to the model is limited in some way (e.g., to registered users), but it should be possible for other researchers to have some path to reproducing or verifying the results.
        \end{enumerate}
    \end{itemize}


\item {\bf Open access to data and code}
    \item[] Question: Does the paper provide open access to the data and code, with sufficient instructions to faithfully reproduce the main experimental results, as described in supplemental material?
    \item[] Answer: \answerTODO{} % Replace by \answerYes{}, \answerNo{}, or \answerNA{}.
    \item[] Justification: \justificationTODO{}
    \item[] Guidelines:
    \begin{itemize}
        \item The answer \answerNA{} means that paper does not include experiments requiring code.
        \item Please see the NeurIPS code and data submission guidelines (\url{https://neurips.cc/public/guides/CodeSubmissionPolicy}) for more details.
        \item While we encourage the release of code and data, we understand that this might not be possible, so \answerNo{} is an acceptable answer. Papers cannot be rejected simply for not including code, unless this is central to the contribution (e.g., for a new open-source benchmark).
        \item The instructions should contain the exact command and environment needed to run to reproduce the results. See the NeurIPS code and data submission guidelines (\url{https://neurips.cc/public/guides/CodeSubmissionPolicy}) for more details.
        \item The authors should provide instructions on data access and preparation, including how to access the raw data, preprocessed data, intermediate data, and generated data, etc.
        \item The authors should provide scripts to reproduce all experimental results for the new proposed method and baselines. If only a subset of experiments are reproducible, they should state which ones are omitted from the script and why.
        \item At submission time, to preserve anonymity, the authors should release anonymized versions (if applicable).
        \item Providing as much information as possible in supplemental material (appended to the paper) is recommended, but including URLs to data and code is permitted.
    \end{itemize}


\item {\bf Experimental setting/details}
    \item[] Question: Does the paper specify all the training and test details (e.g., data splits, hyperparameters, how they were chosen, type of optimizer) necessary to understand the results?
    \item[] Answer: \answerTODO{} % Replace by \answerYes{}, \answerNo{}, or \answerNA{}.
    \item[] Justification: \justificationTODO{}
    \item[] Guidelines:
    \begin{itemize}
        \item The answer \answerNA{} means that the paper does not include experiments.
        \item The experimental setting should be presented in the core of the paper to a level of detail that is necessary to appreciate the results and make sense of them.
        \item The full details can be provided either with the code, in appendix, or as supplemental material.
    \end{itemize}

\item {\bf Experiment statistical significance}
    \item[] Question: Does the paper report error bars suitably and correctly defined or other appropriate information about the statistical significance of the experiments?
    \item[] Answer: \answerTODO{} % Replace by \answerYes{}, \answerNo{}, or \answerNA{}.
    \item[] Justification: \justificationTODO{}
    \item[] Guidelines:
    \begin{itemize}
        \item The answer \answerNA{} means that the paper does not include experiments.
        \item The authors should answer \answerYes{} if the results are accompanied by error bars, confidence intervals, or statistical significance tests, at least for the experiments that support the main claims of the paper.
        \item The factors of variability that the error bars are capturing should be clearly stated (for example, train/test split, initialization, random drawing of some parameter, or overall run with given experimental conditions).
        \item The method for calculating the error bars should be explained (closed form formula, call to a library function, bootstrap, etc.)
        \item The assumptions made should be given (e.g., Normally distributed errors).
        \item It should be clear whether the error bar is the standard deviation or the standard error of the mean.
        \item It is OK to report 1-sigma error bars, but one should state it. The authors should preferably report a 2-sigma error bar than state that they have a 96\% CI, if the hypothesis of Normality of errors is not verified.
        \item For asymmetric distributions, the authors should be careful not to show in tables or figures symmetric error bars that would yield results that are out of range (e.g., negative error rates).
        \item If error bars are reported in tables or plots, the authors should explain in the text how they were calculated and reference the corresponding figures or tables in the text.
    \end{itemize}

\item {\bf Experiments compute resources}
    \item[] Question: For each experiment, does the paper provide sufficient information on the computer resources (type of compute workers, memory, time of execution) needed to reproduce the experiments?
    \item[] Answer: \answerTODO{} % Replace by \answerYes{}, \answerNo{}, or \answerNA{}.
    \item[] Justification: \justificationTODO{}
    \item[] Guidelines:
    \begin{itemize}
        \item The answer \answerNA{} means that the paper does not include experiments.
        \item The paper should indicate the type of compute workers CPU or GPU, internal cluster, or cloud provider, including relevant memory and storage.
        \item The paper should provide the amount of compute required for each of the individual experimental runs as well as estimate the total compute. 
        \item The paper should disclose whether the full research project required more compute than the experiments reported in the paper (e.g., preliminary or failed experiments that didn't make it into the paper). 
    \end{itemize}
    
\item {\bf Code of ethics}
    \item[] Question: Does the research conducted in the paper conform, in every respect, with the NeurIPS Code of Ethics \url{https://neurips.cc/public/EthicsGuidelines}?
    \item[] Answer: \answerTODO{} % Replace by \answerYes{}, \answerNo{}, or \answerNA{}.
    \item[] Justification: \justificationTODO{}
    \item[] Guidelines:
    \begin{itemize}
        \item The answer \answerNA{} means that the authors have not reviewed the NeurIPS Code of Ethics.
        \item If the authors answer \answerNo, they should explain the special circumstances that require a deviation from the Code of Ethics.
        \item The authors should make sure to preserve anonymity (e.g., if there is a special consideration due to laws or regulations in their jurisdiction).
    \end{itemize}


\item {\bf Broader impacts}
    \item[] Question: Does the paper discuss both potential positive societal impacts and negative societal impacts of the work performed?
    \item[] Answer: \answerTODO{} % Replace by \answerYes{}, \answerNo{}, or \answerNA{}.
    \item[] Justification: \justificationTODO{}
    \item[] Guidelines:
    \begin{itemize}
        \item The answer \answerNA{} means that there is no societal impact of the work performed.
        \item If the authors answer \answerNA{} or \answerNo, they should explain why their work has no societal impact or why the paper does not address societal impact.
        \item Examples of negative societal impacts include potential malicious or unintended uses (e.g., disinformation, generating fake profiles, surveillance), fairness considerations (e.g., deployment of technologies that could make decisions that unfairly impact specific groups), privacy considerations, and security considerations.
        \item The conference expects that many papers will be foundational research and not tied to particular applications, let alone deployments. However, if there is a direct path to any negative applications, the authors should point it out. For example, it is legitimate to point out that an improvement in the quality of generative models could be used to generate Deepfakes for disinformation. On the other hand, it is not needed to point out that a generic algorithm for optimizing neural networks could enable people to train models that generate Deepfakes faster.
        \item The authors should consider possible harms that could arise when the technology is being used as intended and functioning correctly, harms that could arise when the technology is being used as intended but gives incorrect results, and harms following from (intentional or unintentional) misuse of the technology.
        \item If there are negative societal impacts, the authors could also discuss possible mitigation strategies (e.g., gated release of models, providing defenses in addition to attacks, mechanisms for monitoring misuse, mechanisms to monitor how a system learns from feedback over time, improving the efficiency and accessibility of ML).
    \end{itemize}
    
\item {\bf Safeguards}
    \item[] Question: Does the paper describe safeguards that have been put in place for responsible release of data or models that have a high risk for misuse (e.g., pre-trained language models, image generators, or scraped datasets)?
    \item[] Answer: \answerTODO{} % Replace by \answerYes{}, \answerNo{}, or \answerNA{}.
    \item[] Justification: \justificationTODO{}
    \item[] Guidelines:
    \begin{itemize}
        \item The answer \answerNA{} means that the paper poses no such risks.
        \item Released models that have a high risk for misuse or dual-use should be released with necessary safeguards to allow for controlled use of the model, for example by requiring that users adhere to usage guidelines or restrictions to access the model or implementing safety filters. 
        \item Datasets that have been scraped from the Internet could pose safety risks. The authors should describe how they avoided releasing unsafe images.
        \item We recognize that providing effective safeguards is challenging, and many papers do not require this, but we encourage authors to take this into account and make a best faith effort.
    \end{itemize}

\item {\bf Licenses for existing assets}
    \item[] Question: Are the creators or original owners of assets (e.g., code, data, models), used in the paper, properly credited and are the license and terms of use explicitly mentioned and properly respected?
    \item[] Answer: \answerTODO{} % Replace by \answerYes{}, \answerNo{}, or \answerNA{}.
    \item[] Justification: \justificationTODO{}
    \item[] Guidelines:
    \begin{itemize}
        \item The answer \answerNA{} means that the paper does not use existing assets.
        \item The authors should cite the original paper that produced the code package or dataset.
        \item The authors should state which version of the asset is used and, if possible, include a URL.
        \item The name of the license (e.g., CC-BY 4.0) should be included for each asset.
        \item For scraped data from a particular source (e.g., website), the copyright and terms of service of that source should be provided.
        \item If assets are released, the license, copyright information, and terms of use in the package should be provided. For popular datasets, \url{paperswithcode.com/datasets} has curated licenses for some datasets. Their licensing guide can help determine the license of a dataset.
        \item For existing datasets that are re-packaged, both the original license and the license of the derived asset (if it has changed) should be provided.
        \item If this information is not available online, the authors are encouraged to reach out to the asset's creators.
    \end{itemize}

\item {\bf New assets}
    \item[] Question: Are new assets introduced in the paper well documented and is the documentation provided alongside the assets?
    \item[] Answer: \answerTODO{} % Replace by \answerYes{}, \answerNo{}, or \answerNA{}.
    \item[] Justification: \justificationTODO{}
    \item[] Guidelines:
    \begin{itemize}
        \item The answer \answerNA{} means that the paper does not release new assets.
        \item Researchers should communicate the details of the dataset\slash code\slash model as part of their submissions via structured templates. This includes details about training, license, limitations, etc. 
        \item The paper should discuss whether and how consent was obtained from people whose asset is used.
        \item At submission time, remember to anonymize your assets (if applicable). You can either create an anonymized URL or include an anonymized zip file.
    \end{itemize}

\item {\bf Crowdsourcing and research with human subjects}
    \item[] Question: For crowdsourcing experiments and research with human subjects, does the paper include the full text of instructions given to participants and screenshots, if applicable, as well as details about compensation (if any)? 
    \item[] Answer: \answerTODO{} % Replace by \answerYes{}, \answerNo{}, or \answerNA{}.
    \item[] Justification: \justificationTODO{}
    \item[] Guidelines:
    \begin{itemize}
        \item The answer \answerNA{} means that the paper does not involve crowdsourcing nor research with human subjects.
        \item Including this information in the supplemental material is fine, but if the main contribution of the paper involves human subjects, then as much detail as possible should be included in the main paper. 
        \item According to the NeurIPS Code of Ethics, workers involved in data collection, curation, or other labor should be paid at least the minimum wage in the country of the data collector. 
    \end{itemize}

\item {\bf Institutional review board (IRB) approvals or equivalent for research with human subjects}
    \item[] Question: Does the paper describe potential risks incurred by study participants, whether such risks were disclosed to the subjects, and whether Institutional Review Board (IRB) approvals (or an equivalent approval/review based on the requirements of your country or institution) were obtained?
    \item[] Answer: \answerTODO{} % Replace by \answerYes{}, \answerNo{}, or \answerNA{}.
    \item[] Justification: \justificationTODO{}
    \item[] Guidelines:
    \begin{itemize}
        \item The answer \answerNA{} means that the paper does not involve crowdsourcing nor research with human subjects.
        \item Depending on the country in which research is conducted, IRB approval (or equivalent) may be required for any human subjects research. If you obtained IRB approval, you should clearly state this in the paper. 
        \item We recognize that the procedures for this may vary significantly between institutions and locations, and we expect authors to adhere to the NeurIPS Code of Ethics and the guidelines for their institution. 
        \item For initial submissions, do not include any information that would break anonymity (if applicable), such as the institution conducting the review.
    \end{itemize}

\item {\bf Declaration of LLM usage}
    \item[] Question: Does the paper describe the usage of LLMs if it is an important, original, or non-standard component of the core methods in this research? Note that if the LLM is used only for writing, editing, or formatting purposes and does \emph{not} impact the core methodology, scientific rigor, or originality of the research, declaration is not required.
    %this research? 
    \item[] Answer: \answerTODO{} % Replace by \answerYes{}, \answerNo{}, or \answerNA{}.
    \item[] Justification: \justificationTODO{}
    \item[] Guidelines:
    \begin{itemize}
        \item The answer \answerNA{} means that the core method development in this research does not involve LLMs as any important, original, or non-standard components.
        \item Please refer to our LLM policy in the NeurIPS handbook for what should or should not be described.
    \end{itemize}

\end{enumerate}

\end{document}
